Space, the final frontier. Many have dreamed to explore it, see the wonders that hide within the vast cosmos of our universe. It has been millennia since space travel became a reality, no longer a dream, a hope, a blurry image taken by primitive technology from ages long past. Intergalactic travel had become a mundane reality. People going to and from in what could best be described as a tin can hurdling through vast empty swathes at speeds sometimes exceeding light itself. It's hard today to imagine this used to be reserved for the best of the best, elite groups that made money in ways shadier than the dark side of the moon, these were the representatives, ambassadors of our people. It's hard to imagine that simple people like me could one day, simply decide to visit the incomprehensible, mind bending void on the simplest of whims.

I have grown blind to the notion of space travel as something exhilarating. Everyday, I wake up, somewhere new each time, either surrounded by the void or, on rarer occasions, on a planet where cargo finally meets its destination. Neither these places nor the people I meet are constants. They're but a fleeting moment, a small piece of the wider cold universe which cares little for us. What I have learned is that the only constants here, are me and the hunk of junk I lovingly call the `Rust Wyrm'. A sleek, long, outdated freighter, with an orange-red body from when that was all the rage, adorned with plates of metal and plastic of many colours, never that same orange-red, each telling a story than no one else could ever truly appreciate. It never made much sense to properly repair the hull in a cohesive manner; faster than light (FTL) travel is costly in many ways, especially on the frail exterior of ships, even modern ones. Despite its age, this bucket of bolts was comfortable, to some extent, and it could still travel efficiently and do its job. 

The interior, much like me and the hull, has seen better days, time seems to have drained the colours from the furniture and art that tried to mimic a humble apartment. Aside from the living quarters, the main areas were the cockpit, the cargo room and the engine room. The first of which I would consider both the best and worst of areas this ship could offer. It's where all the action happened that became so romanticised. Different media always depict a bustling crew in this room, with people getting ready for combat or liftoff, with lively, chaotic conversations about how their poor ship, unsuitable for battle, could never stand up against an armada. It's quiet, the last war was generations ago, the visuals of explosions - replaced by a slab of black granite in which all life seemingly exists. The wide bay window gives into a captain's, or more aptly, a driver's chair, well worn in, next to a console with faded buttons who's inscriptions became useless with time. Every time I sit there, I get to enjoy the view of something I have seen again and again, combined with the reminder that I get to argue with myself over the fate of this ship as I outrun an army of delivery deadlines with ease.

The cargo and engine rooms are both important and insignificant. Most times, I am only there to make sure the shipments are still here, make sure the engine has yet to give out, or simply for a change of scenery. Recently though, something strange has happened, the colour of these rooms had been fading as well. Even more unusual is that no one has ever complained about it, the load delivered was always accepted, I've yet to hear someone complain that their shipment of furniture or what else was discoloured. I thought this was FTL sickness, a side effect of travelling too fast for too long, but this isn't one of them. Usually it's nausea, with rare cases of hallucinations, especially when travelling alone. Something was happening and I did not know what, but it mattered not. I still had a job to do. People, time, the universe, even I, cared little for what a small piece of the universe felt.

In order to achieve FLT speed for a consistent amount of time, ships make use of worm-holes. Some use naturally available ones, but better equipped ships can create their own. Depending on the destination and the sum of the payout, I make a usage of either. Despite the banal nature of going from one place to another, every single day, these gateways through time and space make it all ever so slightly worth it. Entering a worm-hole is something few still get to experience despite the prevalence of FTL travel, as it's usually short lived for the commute most need to make. To enter such a tunnel is to experience the impossible. Inside, you get to see alternate realities, images of worlds with sometimes only the slightest of variations. In the infinity of the universe, some things simply don't happen, and those moments of travel let you see what could have been but is too late to change.

This is possibly the last flight I will ever make, I simply cannot go on. Colours no longer drain as even the exterior of the Wyrm, with all its dents and haphazardly repairs, stares at me in monotone black and white, the same one I've grown sick and tired from seeing in a void that neither answers nor stares back. One last jump to end it all, one last jump to put the final nail in the black hole that is this cursed career, one last jump before the `Rust Wyrm' gets sold as a the piece of scrap it is. And so, I enter the passageway that would lead to the resting place of what I carry this time. Lights of different fantastical colours flood the room, fantastical black, grey, white. As always, I start the hear the creaks and noises of plates being torn off by the journey-. I try and take in the scenery and absorb the mysteries that the multiverse hides. Planets that I've visited, in one world never existed or cessed to be. More plates are violently pulled away. In another reality, some person I met became a wealthy business owner. Red lights shine from the console, replacing the usual light blue glow-. And in yet another one, space travel never left the realm of the imaginary. Alarms start sounding-. In one world, I never traveled as I realized much sooner that the wealth made from this venture, was never worth it. The computer whines and screams about the reactor, the hull, the life support-. Never has one of these images truly captivated me, they were interesting in their own right, but what use did I have for impossibilities that involved those I didn't know. Systems panic from the overwhelming stress-. At least, not until now. The ship was never meant to be in FTL this long-. In some ridiculous twist of fate, the cosmos has given me an answer to a question I never dared to ask the very same moment I answered it. The `Rusty Wyrm' begs and pleads-. Was success ever worth the loss of self ? Am I really more than a small piece the universe can toss away ? The large window cracks, the cargo is engulfed in flames, the engine has given up, the cacophony of alarms turns to eery silence from a simple push of a button that looks brand new-. I stare at this image of me, blinded by how bright and colourful it is. In the final few moments, the cold universe is replaced by the warm embrace of the `Rust Wyrm'-.